\chapter{잔차 스트림, 회로와 내부 알고리즘}
\label{chap:phase5}

\section{결론부터: 회로는 행동에 상대적인 계산 부분그래프다}

회로는 명시된 입력 분포와 행동 metric에 대해 모델의 계산을 충분히 재현하거나 설명한다고
가정한 구성요소와 연결의 부분그래프다. Transformer 전체도 하나의 계산그래프지만, 모든
edge를 회로라고 부르면 설명 단위가 너무 커진다. 따라서 ``모델의 회로''라는 표현보다 다음
tuple이 정확하다.
\index{circuit@회로(circuit)}
\begin{equation}
  \mathcal{C}=(\text{model},\text{input distribution},
  \text{behavior},\text{nodes},\text{edges},\text{intervention rule}).
\end{equation}

같은 모델도 이름 복사, 사실 회상과 거절에는 다른 회로 가설을 가질 수 있다. 같은 행동도
prompt 형식이나 metric이 바뀌면 중요한 경로가 달라질 수 있다. 회로가 발견되었다는 말은
이 tuple의 범위에서 평가해야 한다.

\begin{limitbox}
``회로''에는 분야 전체가 합의한 유일한 원자 단위가 없다. node를 attention head, MLP
block, neuron, SAE feature 또는 token-position별 activation으로 잡을 수 있다. 서로 다른
분해의 회로 크기와 완전성을 그대로 비교해서는 안 된다.
\end{limitbox}

\section{residual stream을 통신 채널로 보기}

pre-norm Transformer에서 각 sublayer는 residual stream의 정규화된 사본을 읽고
$d$차원의 update를 다시 쓴다. 단순화하면 마지막 residual state는
\begin{equation}
  \vect{x}^{(L)}_t=\vect{x}^{(0)}_t+
  \sum_{\ell=0}^{L-1}\Delta\vect{x}^{(\ell)}_{\mathrm{attn},t}+
  \sum_{\ell=0}^{L-1}\Delta\vect{x}^{(\ell)}_{\mathrm{mlp},t}
  \label{eq:residual-sum}
\end{equation}
로 정확히 쓸 수 있다. 각 update의 값은 앞선 모든 update를 normalization과 비선형
연산으로 읽은 결과이므로 서로 독립적인 항은 아니다.

Elhage 등은 이 구조를 여러 구성요소가 공용 residual stream을 통해 읽고 쓰는 통신
채널로 해석하고, attention-only 소형 모델에서 경로 전개와 QK/OV 분해를 체계화했다
\citep{elhage2021framework}. 이 관점은 architecture의 덧셈 구조와 잘 맞지만, residual
vector의 의미가 자연스럽게 분리된 slot으로 보장된다는 뜻은 아니다. 서로 다른 update가
같은 subspace를 사용하거나 상쇄할 수 있다.

\begin{figure}[htbp]
\centering
\begin{tikzpicture}[
  node distance=7mm and 5mm,
  box/.style={draw=Navy,rounded corners,fill=SoftBlue,minimum width=22mm,
              text width=22mm,minimum height=8mm,align=center,font=\small},
  arr/.style={-{Latex[length=2mm]},thick,Navy}
]
  \node[box] (r0) {residual\\$\vect{x}^{(\ell)}$};
  \node[box,right=of r0] (a) {attention\\update};
  \node[box,right=of a] (r1) {residual\\$\vect{u}^{(\ell)}$};
  \node[box,right=of r1] (m) {MLP\\update};
  \node[box,right=of m] (r2) {residual\\$\vect{x}^{(\ell+1)}$};
  \draw[arr] (r0) -- (a);
  \draw[arr] (a) -- node[above,font=\scriptsize] {$+$} (r1);
  \draw[arr] (r1) -- (m);
  \draw[arr] (m) -- node[above,font=\scriptsize] {$+$} (r2);
  \draw[arr] (r0.south) to[out=-45,in=-135] node[below,font=\scriptsize] {identity path} (r1.south);
  \draw[arr] (r1.south) to[out=-45,in=-135] node[below,font=\scriptsize] {identity path} (r2.south);
\end{tikzpicture}
\caption{한 block에서 residual stream을 읽고 update를 더하는 경로}
\label{fig:residual-communication}
\end{figure}

\section{direct logit attribution과 경로 전개}

DLA는 final normalization을 선형화하거나 생략한 조건에서 residual update와
unembedding의 내적으로 직접 기여를 배정한다. 이 조건에서 token $v$의 logit은
\begin{equation}
  z_v=\vect{w}_{U,v}^{\mathsf T}\vect{x}^{(L)}_t+b_v
\end{equation}
다. 
\Cref{eq:residual-sum}을 대입하면 embedding과 각 sublayer update가 주는 직접 내적을
계산할 수 있다. 이를 direct logit attribution(DLA)이라고 부른다. 비교 대상 token
$v^+$와 $v^-$의 logit difference를 쓰면 구성요소 $c$의 직접 기여는
\begin{equation}
  \operatorname{DLA}_c=
  (\vect{w}_{U,v^+}-\vect{w}_{U,v^-})^{\mathsf T}\Delta\vect{x}_c
  \label{eq:dla}
\end{equation}
로 둘 수 있다.

DLA는 residual에 최종적으로 쓰인 방향이 unembedding과 얼마나 정렬되는지 측정한다.
구성요소 $c$가 이후 layer의 query, key 또는 MLP gate를 바꾸는 간접 효과는
\cref{eq:dla}에 포함되지 않는다. final LayerNorm이나 RMSNorm도 모든 성분을 결합하므로
정확한 attribution 규칙을 명시해야 한다.

경로 전개는 residual connection을 따라 계산을 합성하여 입력에서 logit까지의 경로를
열거하는 대수적 관점이다. attention pattern과 normalization scale을 특정 입력에서
고정하면 attention의 value-output 경로는 선형이므로 matrix product로 합칠 수 있다.
그러나 MLP activation, attention softmax와 normalization scale은 원래 입력 의존적이다.
이를 고정한 전개는 해당 입력 주변의 조건부 설명이지 전역 선형 모델이 아니다.

\section{QK 회로와 OV 회로}

QK 회로는 destination 표현과 source 표현 사이의 attention score를 정하는 bilinear
form이다. column-vector convention으로 정규화된 residual을 $\vect{r}_i$라고 두고
\begin{equation}
  \vect{q}_i=\mat{W}_Q^{\mathsf T}\vect{r}_i,\qquad
  \vect{k}_j=\mat{W}_K^{\mathsf T}\vect{r}_j,\qquad
  \vect{v}_j=\mat{W}_V^{\mathsf T}\vect{r}_j
\end{equation}
라고 하자. score의 핵심 bilinear form은
\begin{equation}
  \vect{r}_i^{\mathsf T}
  \underbrace{\mat{W}_Q\mat{W}_K^{\mathsf T}}_{\text{QK matrix}}
  \vect{r}_j
\end{equation}
다. 이 form은 어떤 destination 표현이 어떤 source 표현을 선택하는지 결정한다.

OV map은 선택된 source의 value가 residual stream에 쓰는 내용을 정한다. 그 update는
\begin{equation}
  \Delta\vect{x}_i=\sum_j A_{ij}
  \underbrace{\mat{W}_O^{\mathsf T}\mat{W}_V^{\mathsf T}}_{\text{OV map}}
  \vect{r}_j
\end{equation}
로 나타낼 수 있다. 정확한 전치 방향은 코드의 row/column convention에 따라 달라지지만,
선택을 정하는 QK와 전달 내용을 정하는 OV를 분리한다는 요점은 같다
\citep{elhage2021framework}.
\index{QK circuit@QK 회로}
\index{OV circuit@OV 회로}

큰 attention weight는 QK 선택에 관한 관찰이다. OV map이 source 정보를 target logit과
반대되는 방향으로 쓰거나 거의 0으로 보내면 출력 기여는 작거나 음수일 수 있다. 반대로
작은 weight도 매우 큰 value-output vector와 결합되면 무시할 수 없을 수 있다.

\section{head composition}

앞 layer의 head가 residual에 쓴 vector는 뒤 head의 query, key와 value projection이
모두 읽을 수 있다. 이에 따라 세 가지 대표 합성이 생긴다.

\begin{description}[style=nextline,leftmargin=2em]
  \item[Q-composition]
  앞 head의 출력이 뒤 head의 query를 바꾸어 destination이 찾는 source를 조절한다.

  \item[K-composition]
  앞 head의 출력이 뒤 head의 key를 바꾸어 특정 source가 선택될 가능성을 조절한다.

  \item[V-composition]
  앞 head가 한 position에 쓴 내용을 뒤 head가 value 경로로 다른 position에 운반한다.
\end{description}

이 분류는 attention-only 모델의 virtual weight를 분석하는 데 유용하다. 실제 깊은 모델에서는
normalization, MLP update, 여러 head의 합과 softmax가 사이에 있으므로 matrix product의
크기만으로 인과 효과를 확정할 수 없다. composition score는 후보 경로를 찾는 도구이며
개입 검증이 뒤따라야 한다.

\section{induction head}

induction head가 구현하는 대표 패턴은
\begin{equation}
  [A][B]\;\cdots\;[A]\longrightarrow[B]
\end{equation}
다. 앞에서 $A$ 다음에 $B$가 나타났다면, 뒤의 $A$ 위치에서 $B$를 다음 token으로
복사하는 경향이다. 두 layer attention-only 모델의 이상화된 회로에서는 다음과 같이
설명할 수 있다.
\index{induction head@induction head}

\begin{enumerate}
  \item 앞 layer의 previous-token head가 각 position에 바로 앞 token의 정보를 쓴다.
  \item 뒤 layer induction head의 query가 현재 token $A$를 나타내고, key가 ``바로 앞
        token이 $A$였던 source''와 맞는다.
  \item 그 source의 value에는 $A$ 다음 token인 $B$의 정보가 있으므로 OV 경로가 $B$
        logit을 높이는 방향으로 이를 복사한다.
\end{enumerate}

Olsson 등은 소형 attention-only 모델에서 induction head와 in-context loss 개선 사이에
강한 인과 증거를 제시하고, MLP가 있는 더 큰 모델에서는 여러 상관적 증거를 보고했다
\citep{olsson2022induction}. 원 논문 자체가 증거 수준을 구분한다. 따라서 ``모든 대형
LLM의 in-context learning은 induction head 하나로 완전히 설명된다''고 요약하면 원
결과의 범위를 넘는다.

\begin{resultbox}{Olsson et al. (2022)의 제한된 결과}
반복 token pattern, training 중 induction-like head의 출현 시점과 in-context learning
metric을 연결한 결과다. 소형 attention-only 모델에 대한 강한 인과 증거와 더 큰 모델에
대한 간접 증거를 같은 수준으로 취급하지 않는다.
\end{resultbox}

\section{IOI 회로}

간접 목적어 식별(IOI)은 ``When Mary and John went to the store, John gave a drink to''와
같은 문맥에서 반복되지 않은 이름인 Mary를 예측하는 과제다. Wang 등은 GPT-2 small에서
이 행동을 분석하여 26개 attention head를 7개 주요 기능군으로 묶은 회로를 제시했다.
연구는 name mover, negative name mover, duplicate-token, previous-token, induction 및
S-inhibition 기능을 포함하고, causal intervention과 faithfulness, completeness,
minimality 기준으로 설명을 평가했다 \citep{wang2022ioi}.
\index{IOI@간접 목적어 식별(IOI)}

이 연구의 중요성은 자연어 행동 하나를 입력에서 logit까지 비교적 세밀한 head-level
경로로 연결했다는 데 있다. 동시에 다음 범위가 명시되어야 한다.

\begin{itemize}
  \item 분석 대상은 GPT-2 small과 연구가 구성한 IOI distribution이다.
  \item node granularity는 주로 attention head이며 MLP와 더 미세한 feature를 완전히
        의미론적으로 분해하지 않는다.
  \item 정량 기준이 설명을 지지했지만 저자들은 남은 설명 격차도 보고했다.
  \item 다른 이름 분포, 문장 구조, 모델과 과제에 동일한 회로가 자동으로 적용되지는 않는다.
\end{itemize}

따라서 IOI는 ``자연 모델도 부분적으로 reverse engineering할 수 있다''는 강한 사례지만,
``자연어 모델 전체의 알고리즘을 발견했다''는 사례는 아니다.

\section{MLP feature와 memory 해석}

기본 FFN은 중간 neuron별 activation과 residual 출력 방향의 합으로 전개할 수 있다.
\begin{equation}
  \operatorname{FFN}(\vect{x})=
  \sum_{i=1}^{d_{\mathrm{ff}}}
  \phi(\vect{k}_i^{\mathsf T}\vect{x}+b_i)\vect{v}_i+\vect{b}_{\mathrm{out}}
  \label{eq:mlp-sum}
\end{equation}
input matrix의 row $\vect{k}_i$는 activation 조건을 정하고, output
matrix의 대응 column $\vect{v}_i$는 residual에 쓰는 방향을 정한다. 이 대수 때문에 한
neuron을 key-value pair처럼 분석할 수 있다.

Geva 등은 Transformer FFN의 input pattern과 output vocabulary 분포를 조사하고, lower
layer에는 상대적으로 얕은 pattern이, upper layer에는 더 semantic한 pattern이 나타나는
경향을 보고하며 ``key-value memory'' 해석을 제안했다 \citep{geva2021ffn}. 이 결과는
\cref{eq:mlp-sum}의 분해와 실험 관찰을 결합한 해석이다.

SwiGLU에서는 두 projection의 곱과 SiLU가 있으므로 neuron 하나를
\cref{eq:mlp-sum}의 단순 key-value pair로 그대로 대응시키기 어렵다. 또한 superposition과
다의성 때문에 neuron보다 feature dictionary가 더 적합할 수 있지만, SAE 역시 근사
분해다. ``MLP는 데이터베이스다''는 비유를 literal한 주소-값 저장 구조로 받아들여서는
안 된다.

\section{완전한 알고리즘을 아는 특수 사례}

Tracr는 제한된 배열 처리 언어로 쓴 human-readable program을 표준 decoder-only
Transformer weight로 컴파일한다. token 빈도 계산, 정렬과 괄호 검사 같은 program을
구성하며, compiler가 weight를 만들었기 때문에 source program과 모델 구조의 대응을
ground truth로 사용할 수 있다 \citep{lindner2023tracr}.
\index{Tracr@Tracr}

Tracr 모델에서는 ``어떤 알고리즘을 구현하는가''를 설계 과정에서 안다. 그러나 자연어
corpus의 gradient descent로 학습한 모델에 Tracr의 알려진 program이 그대로 존재한다는 뜻은
아니다. Tracr의 가치는 다음 두 가지다.

\begin{enumerate}
  \item 해석 방법이 알려진 component와 edge를 복원하는지 평가할 수 있다.
  \item 알고리즘을 Transformer 연산으로 구현하는 구체적 construction을 연구할 수 있다.
\end{enumerate}

또 다른 완전성은 모든 floating-point 연산을 기록하는 것이다. 이는 실행 trace로서는
완전하지만, billions of scalar operations를 그대로 나열하므로 사람이 이해하는 압축된
알고리즘 설명과 다르다. 자연 학습 LLM에서 요구하는 어려운 목표는 실행 trace와 행동
사이의 충실하고 압축된 중간 설명을 찾는 일이다.

\begin{factbox}
알려진 program을 컴파일한 Transformer에는 완전한 설계 설명이 존재한다. 학습된 소형
모델에는 특정 과제의 상당한 회로 설명이 존재한다. 자연 학습한 frontier-scale LLM
전체에는 아직 이에 상응하는 전역적이고 완전한 설명이 없다.
\end{factbox}

\section{Phase 5 학습 점검}

\begin{studybox}
\begin{enumerate}
  \item attention-only 2-layer 모델에서 embedding에서 unembedding까지 가능한 residual
        path를 전개한다.
  \item attention head 하나의 QK bilinear form과 OV map을 계산하고, attention weight와
        direct logit contribution을 따로 시각화한다.
  \item 반복 token dataset으로 previous-token head와 induction head 후보를 찾은 뒤
        각각을 ablate하여 loss 변화를 측정한다.
  \item IOI 논문의 node, edge, 입력 분포, logit-difference metric과 개입 기준을 한 장의
        회로 명세로 다시 작성한다.
  \item Tracr program 하나를 컴파일하고 알려진 source variable과 residual direction의
        대응을 blind analysis 결과와 비교한다.
\end{enumerate}
\end{studybox}

\begin{limitbox}
회로 diagram이 그럴듯하다는 사실은 충분하지 않다. 다음 장에서는 node와 edge를 실제로
바꾸어 회로 가설이 원 모델의 반사실적 행동을 예측하는지 검사한다.
\end{limitbox}
